\section{前向过程详解}

\subsection{前向转移分布的定义}

前向过程的每一步转移被定义为一个特定形式的高斯分布：

$$q(x_t | x_{t-1}) = \mathcal{N}\big(x_t;\, \sqrt{1 - \beta_t}\, x_{t-1},\; \beta_t I\big)$$

其中各符号的含义如下：
\begin{itemize}
    \item $\beta_t \in (0, 1)$：第 $t$ 步的噪声调度参数（noise schedule），是关于时间 $t$ 的\textbf{非递减}函数
    \item $\sqrt{1 - \beta_t}\, x_{t-1}$：均值——将上一步的值向原点方向缩放
    \item $\beta_t I$：方差——注入各向同性的高斯噪声
\end{itemize}

\begin{importantbox}{前向过程的直觉理解}
每一步前向转移做了两件事：（1）将当前值 $x_{t-1}$ 乘以一个略小于 1 的系数 $\sqrt{1-\beta_t}$，使其向原点收缩；（2）添加方差为 $\beta_t$ 的高斯噪声。当 $\beta_t$ 接近 0 时，$x_t \approx x_{t-1}$（几乎不变）；当 $\beta_t$ 接近 1 时，$x_t$ 几乎变为纯噪声。随着 $t$ 增大，$\beta_t$ 逐渐增大，噪声逐步积累，最终 $x_T$ 趋近于标准高斯分布。
\end{importantbox}

\subsection{前向跳跃分布（Forward Jump Distribution）}

前向过程的一个关键优势是：我们不需要逐步执行前向过程，可以直接从 $x_0$ 采样任意时间步 $t$ 的 $x_t$。定义 $\alpha_t = 1 - \beta_t$ 和 $\bar{\alpha}_t = \prod_{s=1}^{t} \alpha_s$，则有：

$$q(x_t | x_0) = \mathcal{N}\big(x_t;\, \sqrt{\bar{\alpha}_t}\, x_0,\; (1 - \bar{\alpha}_t) I\big)$$

各符号说明：
\begin{itemize}
    \item $\bar{\alpha}_t = \prod_{s=1}^{t}(1-\beta_s)$：累积保留系数，随 $t$ 递减
    \item $\sqrt{\bar{\alpha}_t}\, x_0$：均值——原始数据的缩放版本
    \item $(1 - \bar{\alpha}_t) I$：方差——累积噪声的总量
\end{itemize}

\begin{importantbox}{前向跳跃分布的重要性}
前向跳跃分布 $q(x_t|x_0)$ 允许我们在训练时\textbf{直接}从 $x_0$ 采样任意时间步的 $x_t$，无需逐步执行前向过程。这是 DDPM 训练高效性的关键所在。利用重参数化技巧：
$$x_t = \sqrt{\bar{\alpha}_t}\, x_0 + \sqrt{1 - \bar{\alpha}_t}\, \epsilon, \quad \epsilon \sim \mathcal{N}(0, I)$$
\end{importantbox}

这个公式建立了 $x_0$、$x_t$ 和噪声 $\epsilon$ 之间的精确关系，后面推导噪声预测目标时会反复用到。

\subsection{本章小结}

前向过程通过预定义的高斯转移分布逐步向数据注入噪声。由于高斯分布的叠加仍为高斯，我们可以推导出闭式的前向跳跃分布 $q(x_t|x_0)$，从而在训练时直接从 $x_0$ 一步采样 $x_t$，无需迭代。

%% ========== 第三节：ELBO 分解与训练目标 ==========

